\documentclass[11pt]{article}
\usepackage[a4paper,margin=2.5cm]{geometry}
\usepackage{amsmath,amssymb,mathtools}
\usepackage[T1]{fontenc}
\usepackage[utf8]{inputenc}

\title{Technical Note: Exact Octic Operator Hierarchy from Linear to Non-Linear Rotors}
\author{Vincenzo Barone}
\date{2026-03-28}

\begin{document}
\maketitle

\section*{Scope}
This note records the exact operator-level results obtained for the octic
rotational sector. The goal is to isolate the part of the theory that is
already rigorous independently of the unresolved compact formulas for the
vibrational dependence of quartic centrifugal distortion constants.

\section*{1. Linear audit of Paper 1 versus Paper 2}
The symbolic audit of the printed Watson--Aliev source equations and of the
compact Aliev linear formulas gives the following exact conclusions:
\begin{enumerate}
\item the compact channel $V$ is derivable from the printed bare $H_{24}$ block;
\item the compact channel $U$ is not derivable from the printed form of
Paper~1;
\item both $\beta_{\parallel}$ and $\beta_{\perp}$ depend essentially on $U$;
\item therefore the compact formulas of Paper~2 for $\beta_{\parallel}$ and
$\beta_{\perp}$ are not derivable symbolically from the printed equations of
Paper~1;
\item the non-compact equations currently retained in the code coincide with
what is derivable from the printed form of Paper~1.
\end{enumerate}

This isolates the unresolved problem to the compact $U$ channel and removes it
from the octic branch.

\section*{2. Exact octic operator hierarchy}
Let $\mathcal O_8$ be the commuting rotational operator space of total degree
eight in $J_x,J_y,J_z$.

\subsection*{2.1 Asymmetric top}
For an asymmetric top the commuting octic basis has dimension
\[
\dim \mathcal O_8^{\mathrm{asym}} = 15.
\]
The general operator is
\[
\mathcal O_8^{\mathrm{asym}}
=
\sum_{a+b+c=4} c_{2a,2b,2c}\,J_x^{2a}J_y^{2b}J_z^{2c}.
\]

\subsection*{2.2 Symmetric top}
For a symmetric top with axis $a$ the operator space collapses to dimension
\[
\dim \mathcal O_8^{\mathrm{sym}} = 5
\]
with basis
\[
J^8,\qquad J^6J_a^2,\qquad J^4J_a^4,\qquad J^2J_a^6,\qquad J_a^8.
\]
Writing the corresponding coefficients as
\[
L_J^{(4)},\quad L_{J^3K},\quad L_{J^2K^2},\quad L_{JK^3},\quad L_K^{(4)},
\]
the exact projection from the cartesian commuting coefficients is
\[
L_J^{(4)}=c_{080},
\]
\[
L_{J^3K}=-4c_{080}+c_{260},
\]
\[
L_{J^2K^2}=6c_{080}-3c_{260}+c_{440},
\]
\[
L_{JK^3}=-4c_{080}+3c_{260}-2c_{440}+c_{620},
\]
\[
L_K^{(4)}=c_{080}-c_{260}+c_{440}-c_{620}+c_{800}.
\]

The exact cartesian image constraints for the $a$-axis symmetric-top subspace
are
\[
c_{422}=2c_{440},\qquad
c_{242}=3c_{260},\qquad
c_{062}=4c_{080},\qquad
c_{044}=6c_{080},
\]
together with the obvious $b\leftrightarrow c$ equalities
\[
c_{602}=c_{620},\qquad
c_{404}=c_{440},\qquad
c_{224}=c_{242},\qquad
c_{026}=c_{062},\qquad
c_{008}=c_{080}.
\]

\subsection*{2.3 Linear rotor}
In the linear limit only one commuting degree-eight operator survives:
\[
\dim \mathcal O_8^{\mathrm{lin}} = 1.
\]
For axis $a$ this operator is
\[
(J_y^2+J_z^2)^4
=
J_y^8+4J_y^6J_z^2+6J_y^4J_z^4+4J_y^2J_z^6+J_z^8.
\]
Therefore the exact linear constraints on the cartesian coefficients are
\[
c_{800}=c_{620}=c_{602}=c_{440}=c_{422}=c_{404}=c_{260}=c_{242}=c_{224}=c_{206}=0,
\]
\[
c_{080}=L,\qquad
c_{062}=4L,\qquad
c_{044}=6L,\qquad
c_{026}=4L,\qquad
c_{008}=L.
\]
Equivalently,
\[
L = c_{080} = L_J^{(4)}.
\]

\section*{3. Consequence for the next derivation step}
The operator algebra is therefore closed. The remaining theoretical problem is
not the projection machinery but the derivation of the cartesian octic
coefficients $c_{2a,2b,2c}$ from the general Watson--Aliev Hamiltonian.

At the current stage the rigorous target is:
\[
\text{Paper 1} \longrightarrow c_{2a,2b,2c} \longrightarrow
\bigl(L_J^{(4)},L_{J^3K},L_{J^2K^2},L_{JK^3},L_K^{(4)}\bigr)
\longrightarrow L.
\]

\section*{4. Data requirements}
For this octic branch there is currently no evidence that third derivatives of
the inverse inertia tensor are needed. The natural tensor hierarchy is instead
the same one already identified for the quartic rotational tensor:
\[
\mu_1,\qquad \mu_2,\qquad \Phi_3.
\]
In particular, $\mu_3$ is not required by the operator analysis developed so
far.

\end{document}
